\section{Ordered endpoints of a whole cell side}
\label{sec:area_law_side_endpoints}

The four sides of a dyadic square are traversed counterclockwise, starting
with the right side. The resulting endpoint identities hold at every
natural scale and every integer cell index; see
\cite[Section~11]{OpenAI2026AreaLaw}.

% Source: 10-geometry.tex, lines 299--306, prop:two-families,
% revision adc7f1241b42e322a6451854ab7e4b4c146bf78a.

\begin{theorem}[Whole-side coordinates and corner witnesses]
    \label{thm:area_law_unsplit_side_endpoints}
    \lean{TNLean.PEPS.AreaLaw.Geometry.cellFan_unsplit_endpoints_coordinates,
        TNLean.PEPS.AreaLaw.Geometry.cellFan_unsplit_endpoints_are_corners}
    \leanok
    \uses{def:area_law_cell_fan,def:area_law_cell_corner}
    Fix an origin $o\in\mathbb R^2$, natural exponent $\ell$ and
    integer cell index $z\in\mathbb Z^2$. Put $t=2^\ell$ and
    $q=o+tz$. If $(a_s,b_s)$ is the ordered pair of endpoints of
    the unsplit side $s\in\{0,1,2,3\}$, then
    \begin{align}
        (a_0,b_0) & =((q_1+t,q_2),(q_1+t,q_2+t)), \\
        (a_1,b_1) & =((q_1+t,q_2+t),(q_1,q_2+t)), \\
        (a_2,b_2) & =((q_1,q_2+t),(q_1,q_2)),     \\
        (a_3,b_3) & =((q_1,q_2),(q_1+t,q_2)).
    \end{align}
    In particular, for every $s$ there exist
    $\varepsilon_s,\eta_s\in\{0,1\}^2$ such that
    \begin{align}
        a_s=o+t(z+\varepsilon_s),\qquad
        b_s=o+t(z+\eta_s).
    \end{align}
    Thus both endpoints are actual corners of the same dyadic cell.
\end{theorem}
\begin{proof}\leanok
    \uses{def:area_law_cell_fan,def:area_law_cell_corner}
    Substitute the four side indices into the affine parametrization
    of the fan boundary. The resulting coordinate pairs are those
    displayed above. The corner pairs $(\varepsilon_s,\eta_s)$,
    in the same order, are
    $((1,0),(1,1))$, $((1,1),(0,1))$,
    $((0,1),(0,0))$ and $((0,0),(1,0))$.
\end{proof}
